% !TEX root = ../nc-quadric-surface-avatar.tex
% 2026-11-23（canonical square完成の相対tails認識を含む）
\section{内在的なモジュライへの忘却}
\label{sec:qa-intrinsic-forgetting}

有限 foundation、周期的 AS 代数、avatar、および構造対象だけを指定した
Abel 圏は、保持する構造も同型の意味も異なる。
本節では、前節までの開部分と内在的モジュライとの関係を、
既に構成した射と、相対的な認識を要する射とに分けて述べる。

\subsection{標的となる強い内在的クラス}

本稿の \(\HH\)、\(\HHstack\) は
\cite[\texttt{def:h2}]{QuadricIntrinsic2026} と
\cite[\texttt{def:relative-h2}]{QuadricModuli2026} の意味で用いる。
特に、任意の三次元 AS 正則 \(\mathbb Z\)-代数の tails 圏を
無条件に \(\HH\) と呼ぶわけではない。

\begin{definition}[体上の強い内在的クラス]
\label{def:qa-h2}
\(\HH\) の対象は \((\mathcal C,\OO)\) であり、次を満たす。
\begin{enumerate}
\item \(\mathcal C\) は \(\kk\)-線形 noetherian Abel 圏で、
\(D^b(\mathcal C)\) は Hom 有限で Serre 関手 \(S\) を持つ。
\item \(\OO\) は rank-one torsion-free exceptional 対象である。
ここで $\mathcal C_0$ は有限長対象の部分圏、$\mathcal C_1$ は
$\mathcal C/\mathcal C_0$ で有限長となる対象の逆像とし、
$q:\mathcal C\to\mathcal C/\mathcal C_1$ とする。
条件は $q\OO$ が非零単純かつ $\Hom(\mathcal C_1,\OO)=0$ である。
また exceptional とは $\End(\OO)=\kk$、$\Ext^{i}(\OO,\OO)=0$ ($i>0$) をいう。
\item 完全自己同値 \(s\) が存在し、
\(s^2\simeq S^{-1}[2]\) かつ \((\OO,s)\) は Artin--Zhang ample である。
\item 同じ \(s\) に対して、全軌道環
\[
 B(\mathcal C,\OO,s)
   =\bigoplus_{n\ge0}\Hom(\OO,s^n\OO)
\]
が \(A/(g)\) と同型である。
ここで \(A\) は左右 noetherian、二次 Koszul AS 正則、
型 \((1,4,6,4,1)\) の環、\(g\in A_2\) は左右正則正規元である。
\end{enumerate}
\(s\)、根同型、\(A\)、\(g\)、商表示は指定構造ではない。
射は \(\OO\) の同型を伴う完全同値であり、
その同型を保つ自然同型によって一射を同一視する。
\end{definition}

ここで型 \((1,4,6,4,1)\) は、自明加群の最小解消の階数列が
この列で、global dimension と Gorenstein parameter がともに四であることを意味する。
\(g\) の中心性は要求しない。
また、ある商圏との同値だけでなく、指定した \(\OO\) と同じ \(s\) の
\emph{全軌道環} に対する表示を要求する。

相対版 \(\HHstack\) は、標数零の代数閉体上の局所有限型スキームの
fppf site 上で定義する。
圏の基底変更は Ind 圏における加群対象の圏によるものであり、
Hom を tensor するだけの基底変更ではない。
対象と完全列の fppf descent、有限 Hom と flat base change、
平坦な \(\OO\) を要求する。
局所 canonical compact 導来圏は相対 smooth proper で、
相対 Serre 双対性と導来基底変更に適合する。
全幾何ファイバーは定義~\ref{def:qa-h2} を満たす。
さらに、\(s\)、AZ 条件、全軌道環の AS4 被覆は
\emph{非被約基底も含めて fppf 局所に存在する}ことを要求する。
根や被覆を重なり上で同一視することは要求しない。
これらは標識ではなく存在条件だからである。

\begin{remark}[相対存在は幾何点条件より強い]
\label{rem:qa-relative-root}
全幾何ファイバーが \(\HH\) に入っても、
相対平方根、AZ ampleness、相対 AS4 被覆が存在することは別問題である。
例外対象の無障害性や有限代数の \(\operatorname{HH}^3\) の消滅だけでは、
これらの条件は得られない。
\cite[\texttt{prop:chosen-root-does-not-lift}]{QuadricModuli2026} は、可換 quadric 上の
ある適合根が、特定の一次 Abel 変形へ持ち上がらない例を与える。
これは別の根の存在まで否定する例ではなく、
\(\HHstack\) の非代数性や許容部分の非開性を証明するものでもない。
\end{remark}

\subsection{canonical square 完成で確立した定義域}

定理~\ref{thm:qa-relative-square-completion-tails} と
命題~\ref{prop:qa-relative-square-unfolding} により，実際の曲線 avatar
から構成した有限 square 代数の canonical 完成と全整数展開には，
標準対象を保つ相対 tails 認識がある．これは幾何点の同値だけではなく，
本稿の site 上の module-object 基底変更と coherent に整合する完全同値である．
左右 Noether 性と係数束付きの実際の AS 解消も証明済みである．
補題~\ref{lem:qa-classified-fiber-finite-square-recognition} は，
古典型を含む分類済み square fiber を持つ有限局所自由族まで扱う．
従ってこの canonical 部分の相対 heart 回復を未解決のままにしない．
一方，任意の別途与えた AS/QP 族との全整数同一視には
注意~\ref{rem:qa-square-unfolding-vertex-lines} の頂点線束・周期比較が必要であり，
以下の一般 presentation stack 全体の認識と区別する．

さらに，定理~\ref{thm:qa-square-normal-foundation-stack-equivalence} は
有限 square と normal foundation の全スタックを同定し，
系~\ref{cor:qa-h2-square-foundation-open} は既存の内在的開をこの表示へ移す．
全 $\HHstack$ については
定理~\ref{thm:qa-universal-h2-rank81-completion} が，ruling の選択を要しない
階数八十一の通常代数と pointed tails を与える．
ここまでの有限表示・canonical 完成の認識は確立済みである．
定理~\ref{thm:qa-square-explicit-h2-open}と
系~\ref{cor:qa-square-h2-twenty-polynomials}は，この全$H_2$開を
三十二次元線形表示上の四組・計二十式以内で記述する．
任意の別途与えた AS/QP presentation 群oidとの比較は後続課題である．

\subsection{忘却射の正確な定義域}

\begin{definition}[内在的条件に適合する族]
\label{def:qa-compatible-chart}
\(\mathfrak U\) を、周期的代数または avatar から
基底変更と同型に関して自然に \((\mathcal C,\OO)\) を構成する
fppf パラメータスタックとする。
\(\mathfrak U_{\HH}\) は、その構成した族が
上記の \(\HHstack\) の相対条件を満たす対象の全部分スタックとする。
この定義において、\(\mathfrak U_{\HH}\) が代数的または開であることは
仮定も結論もしない。
\end{definition}

\begin{proposition}[適合族からの忘却]
\label{prop:qa-compatible-forgetting}
定義~\ref{def:qa-compatible-chart} の構成は射
\[
 \pi_{\mathfrak U}:\mathfrak U_{\HH}\longrightarrow\HHstack,
 \qquad u\longmapsto(\mathcal C_u,\OO_u)
\]
を与える。
二つの chart 間の変換が、基底変更と両立する pointed exact 同値を
与えるなら、二つの忘却射はその変換に沿って自然に二同型となる。
\end{proposition}
\begin{proof}
\(\mathfrak U_{\HH}\) の条件は fppf 局所的で、
pointed exact 同値と基底変更に安定である。
従って同型で飽和した全部分スタックを定める。
構成した圏の系、構造対象、その比較同型を保持して
残る標識を忘れれば、定義された標的への functor を得る。
二つの構成を結ぶ pointed exact 同値が任意の基底変更と整合する場合には、
それらの同値がそのまま二つの functor の二同型を与える。
\end{proof}

この命題は忘却の意味を明示するものであり、
未知の条件を有限個の方程式で認識したという定理ではない。
特に \(\mathfrak F_{\le12}\) は有限代数だけを分類するので、
そこから直接 \(\mathcal C\) を取ることはできない。
\(\operatorname{Perf}(A)\) の構成だけでは、元の曲面の heart は選ばれていない。
有限代数の標準的な加群 heart と、tilting によって移された曲面の heart を
同一視してはならない。

「十二本矢の適当な開部分から \(\HHstack\) への自然な射」という主張を
有限な表示で実現するためには、少なくとも次を示す必要がある。
\begin{enumerate}
\item 使用する periodic algebra、grading、foundation と heart を
相対的に構成し、全次数の平坦性、noetherian 性、基底変更を確認する。
\item 指定した \(\OO\) の相対例外性・rank 条件、相対 Serre 公式、
適合する完全平方根と AZ ampleness を確認する。
\item 同じ根の全軌道環が、相対 AS4 環の正則正規二次切断となることを確認する。
\item これらの条件を持つ部分関手が、考えている表示スタックの
不変開部分で表されることを証明し、表示変更の全ての射を降下させる。
\end{enumerate}
根と全軌道環が既に与えられた場合、固定軌道環への AS4 被覆を選ぶ空間は
\cite[\texttt{thm:relative-finite-as-cover}]{QuadricModuli2026} により、
\(\operatorname{Gr}(6,K)\) 内の有限表示局所閉 incidence として記述できる。
しかし、その被覆を忘れた fppf 像の開性・代数性や、根そのものの存在は
この有限判定の結論ではない。

\subsection{十二本矢の全体を送ることへの具体的な障害}

\begin{lemma}[反 Serre 軌道の生成は必要である]
\label{lem:qa-antiserre-necessary}
\((\mathcal C,\OO)\in\HH\) ならば、
\(t=S^{-1}[2]\) に対する \((\OO,t)\) も AZ ample である。
特に、非零対象 \(M\) に対して
\(\Hom(t^{-n}\OO,M)=0\) が全ての \(n\ge0\) で成立することはない。
\end{lemma}
\begin{proof}
定義の根 \(s\) を選び、\(E=\OO\) と置く。
\(B=A/(g)\) は次数一で生成される。
AZ の生成条件を \(sE\) に適用し、得られた全射に
\(s^{-1}\) を作用させると、\(E\) は
\(s^{-m}E\)（\(m>0\)）の有限直和の商になる。
次数一生成により各成分 \(s^{-m}E\to E\) は、
最後の段が \(s^{-1}E\to E\) である合成の線形結合である。
従って評価 \(\Hom(s^{-1}E,E)\otimes s^{-1}E\to E\) は全射となる。
Hom 有限性により、この評価の源は有限直和である。
これは \cite[\texttt{lem:ample-square-from-generation}]{QuadricIntrinsic2026}
の評価全射である。

任意の対象への AZ 生成全射で奇数指数 \(m\) の項があれば、
この評価全射を \(s^{-m}\) で移して \(s^{-(m+1)}E\) の
有限直和に置き換える。従って偶数指数の項だけで生成できる。
指数に任意の下限を課す場合は、先に対象を正に shift して
生成全射を作り、負の shift で戻す。
全射に対する Hom の最終的な全射性は
元の AZ 条件を偶数部分列に制限すれば従う。
\(s^2\simeq t\) なので \((E,t)\) は AZ ample であり、
最後の主張はその生成条件から従う。
\end{proof}

\begin{example}[重み付き十二本矢と residual gerbe]
\label{ex:qa-weighted-not-h2}
\(S=\kk[x,y,z]\)、\(\deg x=\deg y=1\)、\(\deg z=2\) とし、
\(\mathcal X=\PP(1,1,2)\) を重み付き射影 \emph{stack} とする。
この環は global dimension 三、Gorenstein parameter 四の
noetherian AS 正則環である。
しかし
\[
 (\coh\mathcal X,L)\notin\HH\qquad(L\in\Pic(\mathcal X)).
\]
\end{example}
\begin{proof}
通常の三変数 Koszul 解消に各変数の重みを入れると
\[
 0\longrightarrow S(-4)
 \longrightarrow S(-3)^2\oplus S(-2)
 \longrightarrow S(-1)^2\oplus S(-2)
 \longrightarrow S\longrightarrow\kk\longrightarrow0
\]
を得る。これは AS 条件と Gorenstein parameter を与える。

\((\OO,\OO(1),\OO(2),\OO(3))\) の有限 End 代数を考える。
隣接する \(x_i,y_i\) が六本、距離二の \(z_0,z_1\) が二本であり、
最小関係は
\[
 \begin{aligned}
 x_1y_0-y_1x_0&=0,&x_2y_1-y_2x_1&=0,\\
 z_1x_0-x_2z_0&=0,&z_1y_0-y_2z_0&=0.
 \end{aligned}
\]
従って relation completion は十二本矢である。
前節の四写像表示では \(a,b\) はともに階数三で、核は交代テンソルの直線、
\(c,d\) は全射である。よってこの点は実際に
\(\mathfrak F_{\le12}\) の最小十二本矢 locus に属する。
この四対象が full strong であることも直接確認できる。
重み付き多項式環の局所コホモロジーから
\(H^1(\mathcal X,\OO(n))=0\) であり、
Serre 双対により \(n\ge-3\) では \(H^2(\mathcal X,\OO(n))=0\) である。
これで strong 性が従う。
三変数の sheaf Koszul 完全列は四つの連続する twist から
全ての twist を生成し、有限生成 graded \(S\)-加群の有限自由解消により
これらは \(D^b(\coh\mathcal X)\) 全体を生成する。

一方、重み付き Euler 列から
\(\omega_{\mathcal X}=\OO(-4)\) であり、\(t=-\otimes\OO(4)\) である。
点 \([0:0:1]\) の residual gerbe を
\(i:B\mu_2\hookrightarrow\mathcal X\) とし、
\(\chi\) を \(\mu_2\) の非自明な一次元表現とする。
任意の可逆層 \(L\) に対し、非零 coherent 対象
\(M_L=i_*(i^*L\otimes\chi)\) を取る。
\(\OO(4n)\) の gerbe 上の指標は自明なので、引戻しと直像の随伴により
\[
 \Hom(L(-4n),M_L)=
 H^0(B\mu_2,i^*\OO(4n)\otimes\chi)=0\qquad(n\ge0).
\]
特に標準 helix の四位相すべてで生成条件が破れる。
補題~\ref{lem:qa-antiserre-necessary} に反するので、主張を得る。
\end{proof}

\begin{example}[\(\Sigma_2\) 自身の構造対象]
\label{ex:qa-hirzebruch-not-h2}
\(Y=\Sigma_2\) の負断面を \(C\)、ファイバーを \(f\) とする。
\(C^2=-2\)、\(-K_Y=2C+4f\) なので、
\(\OO_Y(-K_Y)|_C\simeq\OO_C\) である。
任意の可逆層 \(L\) に対し、
\(M_L=i_*(i^*L\otimes\OO_C(-1))\ne0\) とおくと
\[
 \Hom\bigl(L\otimes\OO_Y(nK_Y),M_L\bigr)
   =H^0(C,\OO_C(-1))=0\qquad(n\ge0).
\]
同じ必要条件から \((\coh\Sigma_2,L)\notin\HH\) となる。
これは \emph{avatar の一点 Hilbert 曲面が \(\Sigma_2\) である}ことと、
元の Abel 圏が \(\coh\Sigma_2\) であることを混同してよいという意味ではない。
Hilbert 曲面の型だけから、対応する非可換圏の \(\HH\) 適合性は決まらない。
\end{example}

この二例は、十二本矢の全 locus、AS 正則性、あるいは
full strong foundation の存在だけで忘却の標的が \(\HHstack\) になる、
という主張を否定する。
適切な開部分が存在する可能性を否定するものではなく、
その開部分には構造対象の反 Serre 軌道に対する生成条件も必要であることを示す。

\subsection{標識、射、既存の表示スタック}

\(\HHstack\) への忘却では \(\OO\) を保持する。
従って、foundation の mutation が構造対象を別の対象へ移す場合には、
そのまま同じ pointed 対象の表示変更と扱えない。
QP の mutation が dg 導来同値を与えても、
指定した heart の完全同値と \(\OO\) の同型は別途必要である。
\cite[\texttt{cor:pointed-derived-exact}]{QuadricIntrinsic2026} は体上の \(\HH\) 対象の間の判定であり、
任意の相対 periodic family について同じ結論を自動的に与えるものではない。

同様に、消去可能な二本の矢 \(x,y\) と項 \(yx\) を付加した生の QP には、
\((x,y)\mapsto(\lambda x,\lambda^{-1}y)\) という自己同型がある。
消去後の表示ではこの自己同型は見えなくなる。
従って、生の安定化 QP groupoid の消去を、
射を補正しないままスタック同値とすることはできない。
有限 foundation の頂点スカラーの剛体化、
周期的環の次数スカラーの除去、
構造対象を保つ自己同値全体での同一視は、それぞれ異なる操作である。

現在の内在的研究との関係をまとめると、
\cite[\texttt{prop:foundation-stack}]{QuadricModuli2026} の有限二写像表示 \(\mathfrak F_0\) は、
一方の長さ二の積を可逆にした foundation chart である。

\begin{proposition}[幾何点における十本矢 chart の被覆]
\label{prop:qa-h2-normal-ten-chart}
体上の \(\HH\) の任意の対象は、
\(\mathfrak F_{\le10}^{(a)}\) に属する normal foundation を持つ。
これは幾何点での主張であり、任意の \(\HHstack\) の族における
foundation の相対的な存在・降下を結論するものではない。
\end{proposition}
\begin{proof}
\cite[\texttt{thm:normal-helix}]{QuadricIntrinsic2026} と
\cite[\texttt{prop:ruling-involution-mutation}]{QuadricModuli2026} により、
\((E_0,E_1,E_2,E_3)=(\OO,X,s\OO,Z)\) という foundation があり、
本節の \(b\) は同型、\(c\) は全射である。
\(E_0,E_1,E_2\) は rank-one torsion-free 対象である。
このような対象間の非零射は単射である。
実際、像が低い Gabriel--Krull 段階に入れば標的の torsion-free 性に反し、
従って rank-one 商では非零射、すなわち同型を誘導する。
その核は低い段階に入るが、始域の torsion-free 性によって零である。
従って二つの隣接する非零射の合成は零にならず、
\(\ker a\) は非零の分解可能テンソルを含まない。
\(b\) は同型なので他方の核は零であり、
\(c\) の全射性から定義~\ref{def:qa-foundation-open} の全条件を満たす。
\end{proof}

相対的な ruling 交換を \(\tau\) とすると、
\[
 \mathfrak Q_{\mathrm{fnd}}=[\mathfrak F_0/S_2]
\]
は有限型 Artin スタックであり、体上では
\(\HH(\kk)\) をその full subgroupoid として実現する
（\cite[\texttt{cor:h2-point-groupoid-foundation}]{QuadricModuli2026}）。
この幾何点の結果だけから、\(\mathfrak Q_{\mathrm{fnd}}\) 全体または
その未特定の開部分が \(\HHstack\) と同値であるとはいえない。

相対的には、同論文の反 Serre 認識によって全忠実な射
\[
 \HHstack\hookrightarrow
 \mathfrak P=[X_3/\PGL_9]
\]
がある。
ここで \(X_3\subset\operatorname{Gr}(56,81)\) は、
九生成元・五十六二次関係の環の三次成分が階数四十九となる
scheme-theoretic な局所閉 locus である。
\[
 X_{\mathrm{as}}=X_3\times_{\mathfrak P}\HHstack
\]
は不変 fppf 部分層として定義され、
\(\HHstack\simeq[X_{\mathrm{as}}/\PGL_9]\) が成立する。
\(X_{\mathrm{as}}\) の有限型・滑らか・既約な純次元83のスキームとしての表現可能性は
相対モジュライ論文の定理~\cite[\texttt{thm:global-h2-smooth-finite-type}]{QuadricModuli2026} と
系~\cite[\texttt{cor:global-xas-smooth-scheme}]{QuadricModuli2026} で証明されている。
ただし $X_{\mathrm{as}}\to X_3$ は一致根対で開ではなく、
局所閉性と明示的な偏極認識方程式は未解決である。
従って本稿の八・十・十二本矢の明示的な開部分は、
この相対認識問題への有限表示の候補を与えるが、
それ自体で内在的モジュライの代数性を解決するものではない。
